% example-report.tex — CANONICAL WORKED EXAMPLE of a schema-forge design report.
% This is the structure/style /polish must mirror. It is a *reference*, not a
% real verified design: the "Sunrise Overdrive" content and numbers are
% illustrative, and the figures it \includegraphics (sunrise.pdf, plot-*.pdf)
% are the per-design artifacts /polish renders for the ACTUAL circuit (the
% schematic via `schema_forge.render schematic` -> SVG -> rsvg-convert -> PDF,
% the plots from the verified sim). When authoring design/design-report.tex:
%   * copy this skeleton, keep every section in this order,
%   * swap in the real circuit's title/spec/BOM/measured numbers,
%   * point the figures at design/schematics/<final>.pdf and the sim plots,
%   * use \schematicpage for the landscape schematic (never a raw figure —
%     it bounds the image inside the header/footer),
%   * build with .claude/skills/polish/assets/build-report.sh.
% Every class feature is exercised below: cover setters + \sfmaketitle, the
% blue/green/olive \callout boxes, \sfhead/\sfgrouprow zebra tables, and
% \schematicpage. See schemaforge-report.cls for the full API.
\documentclass{schemaforge-report}

\setdoctitle{Sunrise Overdrive}
\setguidekind{Design \& Build Guide}
\setfootleft{two-transistor silicon drive}
\setfootstatus{verified 10/10 $\cdot$ SPICE}

\setcovereyebrow{SCHEMA-FORGE $\cdot$ VERIFIED BY SIMULATION}
\setcovertitle{Sunrise\\Overdrive}
\setcoversubtitle{A two-transistor silicon overdrive with a symmetric diode clipper}
\setcovertopics{Design rationale $\cdot$ schematic $\cdot$ simulated results $\cdot$ bill of materials $\cdot$ build notes}
\setcoverstatus{ngspice: converged $\cdot$ 10\,/\,10 spec assertions pass $\cdot$ exit 0}
\setcovermeta{%
  \metalabel{Topology} & Two common-emitter gain stages around a symmetric silicon-diode clipper\\
  \metalabel{Supply}   & 9\,V single rail, audio ground\\
  \metalabel{Voicing}  & smooth, mid-forward drive $\approx$11\% THD with a passive treble-cut tone\\
  \metalabel{Status}   & \textbf{10/10 verified} $\cdot$ single converged ngspice run\\
  \metalabel{Generated}& 2026-06-30\\
}

\begin{document}
\sfmaketitle

\begin{callout}{blue}{What you are holding}
This is the complete engineering record for a discrete \textbf{silicon overdrive} --- a smooth,
mid-forward drive built from two common-emitter stages around a symmetric diode clipper --- and
\emph{verified by SPICE simulation} rather than by ear. Every performance claim in this guide is
backed by a machine-checked assertion: the design is called ``verified'' only because ngspice
converged \emph{and} all ten target assertions passed against the measured results. Where a number
appears (a bias point, a gain, a distortion figure), it is the value the simulator measured.
\end{callout}

\section{What this is}
The overdrive is the gentler cousin of the fuzz: instead of letting the transistors themselves
collapse into a gated square wave, a dedicated \emph{clipping element} rounds the peaks off a clean,
amplified signal. Here that element is a pair of back-to-back silicon diodes shunting the signal to
ground. Below their forward drop the signal passes untouched; above it the diodes conduct and softly
flatten the peaks, adding the harmonic warmth that defines the drive.

This guide documents a \textbf{two-stage discrete} version on a single 9\,V rail. The first stage
sets the drive into the clipper; the second recovers level and drives a passive tone and volume
network. Three design goals shape it:
\begin{itemize}
  \item \textbf{Symmetric clipping} for an even, vocal break-up rather than a harsh buzz.
  \item \textbf{A stable bias} so the operating point does not wander with supply or temperature.
  \item \textbf{A buffer-friendly input} and a quiet, mix-ready output stage.
\end{itemize}

\begin{callout}{green}{The honesty contract}
A design was only ever tagged \textbf{verified} when two things were true at once: ngspice
\emph{converged} on a solution (a broken or unphysical circuit will not solve), \emph{and} every
assertion in the target specification passed against the measured \texttt{.measure}/\texttt{.control}
results. No assertion was weakened to make a design pass.
\end{callout}

\section{How it works --- the signal path}
\subsection{The two gain stages}
The guitar enters through a series input resistor and a coupling capacitor onto the base of
\texttt{Q1}, a common-emitter stage biased from a stiff divider. Its emitter degeneration is
AC-bypassed, so the stage runs at full gain while the unbypassed DC path holds the bias steady. The
amplified signal is then AC-coupled into the clipper.

\subsection{The diode clipper}
Two silicon diodes (\texttt{D1}, \texttt{D2}) sit anti-parallel from the clipper node to ground. The
symmetry is what gives this voicing its smooth, even character --- the positive and negative
excursions are flattened identically, so the harmonic content is dominated by odd harmonics with a
gentle roll-off. A second stage, \texttt{Q2}, recovers the level lost across the clipper's series
resistance and presents a low impedance to the tone network.

\begin{callout}{olive}{Why the clipper sits between two stages}
The clipper is deliberately placed \emph{after} the first gain stage and \emph{before} the second.
Clipping a small signal directly off the pickup would need impractically low diode drops; clipping a
fully recovered output would waste headroom. Driving the diodes from a stage already at full gain
puts the signal well past the silicon knee at normal playing level, while the recovery stage restores
a usable output --- the gain budget is split so each block does one job.
\end{callout}

\section{Research \& prior art}
The topology is the canonical ``two-transistor boost into a shunt diode clipper'' that underlies a
whole family of discrete drives. The reference values below were the starting point; the simulator
then tuned them to the spec.
\begin{table}[h]
\centering\small\sffamily
\rowcolors{2}{white}{sfZebra}
\begin{tabular}{@{}lll@{}}
\sfhead{Element} & \sfheadcell{Stock value} & \sfheadcell{Role}\\
Input cap     & 0.1\,$\mu$F & couples the guitar onto \texttt{Q1}\\
\texttt{Q1} load   & 10\,k$\Omega$ & first-stage gain\\
Clipper diodes& 1N4148 $\times$2 & symmetric silicon clipping\\
Clipper load  & 47\,k$\Omega$ & sets drive into the diodes\\
\texttt{Q2} load   & 10\,k$\Omega$ & recovery stage\\
Tone cap      & 10\,nF & passive treble-cut pivot\\
Volume pot    & 100\,k$\Omega$ & master output level\\
\end{tabular}
\caption{Reference discrete-overdrive values used as the starting point.}
\end{table}

\section{The specification}
The target was captured as ten machine-checkable assertions before any circuit was drawn. Each names
a measurement the netlist must produce; the design is \textbf{verified} only when all ten pass in one
converged run. The measured column is the simulator's own parse of the final netlist.

\begin{table}[h]
\centering\small\sffamily
\rowcolors{2}{white}{sfZebra}
\begin{tabular}{@{}lllrll@{}}
\sfhead{Sub-goal} & \sfheadcell{Measure} & \sfheadcell{Target} & \sfheadcell{Measured} & \sfheadcell{Unit} & \sfheadcell{Status}\\
bias\_q1  & vc1\_bias  & 4.0--6.0 (between) & 4.92  & V  & \textbf{\textcolor{sfGreen}{PASS}}\\
bias\_q2  & vc2\_bias  & 4.0--6.0 (between) & 5.08  & V  & \textbf{\textcolor{sfGreen}{PASS}}\\
gain      & gain\_db   & $\geq 30$          & 38.50 & dB & \textbf{\textcolor{sfGreen}{PASS}}\\
drive\_thd& thd\_max   & $\geq 8$           & 11.20 & \% & \textbf{\textcolor{sfGreen}{PASS}}\\
cleanup   & thd\_clean & $\leq 3$           & 0.94  & \% & \textbf{\textcolor{sfGreen}{PASS}}\\
symmetry  & clip\_asym & $\leq 5$           & 1.80  & \% & \textbf{\textcolor{sfGreen}{PASS}}\\
input\_z  & zin        & $\geq 100$         & 138.0 & k$\Omega$ & \textbf{\textcolor{sfGreen}{PASS}}\\
low\_corner& f\_lo     & $\leq 90$          & 38.00 & Hz & \textbf{\textcolor{sfGreen}{PASS}}\\
tone\_range& tone\_db  & $\geq 8$           & 12.40 & dB & \textbf{\textcolor{sfGreen}{PASS}}\\
out\_swing & vout\_pp  & $\geq 2.0$         & 4.40  & V  & \textbf{\textcolor{sfGreen}{PASS}}\\
\end{tabular}
\caption{Target spec versus measured results --- 10/10 verified (ngspice exit 0).}
\end{table}

\section{Schematic}
The schematic on the following page is not a hand-drawing --- it is rendered \emph{directly from the
verified netlist} by a deterministic auto-layout, so it cannot drift from the circuit that was
simulated. Every device, value and connection is exactly what ngspice solved.

\textbf{How to read it.} The audio ground is node \texttt{0}; the supply rail is \texttt{vcc}. Each
net is a horizontal rail and each device a rung between its nets. Trace the signal: \texttt{in}
$\rightarrow$ \texttt{Rin} $\rightarrow$ \texttt{Cin} into \texttt{Q1}'s base; \texttt{Q1}'s
collector drives \texttt{Ccp} into the diode clipper \texttt{(D1/D2)}; the recovered signal leaves
\texttt{Q2}'s collector through \texttt{Cout} into the tone/volume network and out.

\schematicpage{sunrise.pdf}{Complete schematic, auto-rendered from the verified netlist \texttt{sunrise.cir} (27 devices; net-rail layout).}

\section{How it was verified}
The trust root is a single command that runs the netlist through ngspice and checks the measured
results against the spec. It returns one of three honest verdicts: \emph{verified} (converged and all
assertions pass), \emph{converged} (solves, but some assertion is unmet), or \emph{failed}
(non-convergent). Convergence itself is a free correctness gate --- a floating node or a circuit with
no DC path to ground simply will not solve.

\begin{callout}{green}{What ``verified'' guarantees here}
ngspice converged, and all ten assertions in the specification passed against the measured
\texttt{.measure}/\texttt{.control} results, in a single run. The numbers in this guide are that run's
measurements.
\end{callout}

\section{Simulated results}
All plots below are drawn from ngspice output of the verified circuit, at the spec's nominal
conditions: 9\,V supply, a 1\,kHz input, drive at maximum, volume at maximum, tone nominal.

\subsection{Transient response --- the drive}
At a 50\,mV input the clean sine is amplified past the diode knee and the peaks are softly flattened
into the rounded, symmetric clipped wave that is the overdrive. The clipping is gentle rather than
gated --- the body of the waveform survives, which is what keeps the drive musical and touch-sensitive.
\begin{figure}[h]\centering
\includegraphics[width=0.82\linewidth]{plot-transient.pdf}
\caption{Three cycles of steady-state output. The input is shown scaled $\times$40 for shape comparison; the output is the real, unscaled clipped waveform.}
\end{figure}

\subsection{Output spectrum}
The harmonic content of the clipped wave is the drive's tonal signature: a fundamental at 1\,kHz
surrounded by a soft harmonic comb dominated by odd harmonics, corresponding to a measured total
harmonic distortion of \textbf{11.2\%}.
\begin{figure}[h]\centering
\includegraphics[width=0.82\linewidth]{plot-spectrum.pdf}
\caption{Output spectrum (ngspice \texttt{spec}, the same DFT method the verified THD figure uses). THD = 11.2\%.}
\end{figure}

\subsection{Frequency response (Bode)}
The small-signal response shows the $\approx$38.5\,dB midband gain, the low-frequency roll-off set by
the coupling capacitors ($-3$\,dB corner at 38\,Hz, safely below the guitar's low E), and the gentle
high-frequency roll-off of the passive tone network.
\begin{figure}[h]\centering
\includegraphics[width=0.82\linewidth]{plot-bode.pdf}
\caption{Output Bode plot. Midband gain 38.5\,dB; low corner 38\,Hz.}
\end{figure}

\section{Bill of materials \& what each part does}
Every device in the verified netlist, with its function and --- because this is a \emph{build} guide
--- a note on what moving its value does. Sub-blocks are grouped.
\begin{table}[h]
\centering\small\sffamily
\rowcolors{2}{white}{sfZebra}
\begin{tabular}{@{}llp{8.0cm}@{}}
\sfhead{Ref} & \sfheadcell{Value} & \sfheadcell{Function --- and the effect of changing it}\\
\sfgrouprow{3}{Input \& first stage (Q1)}\\
Rin   & 68\,k$\Omega$  & Series input resistor; sets input impedance. \emph{Larger} $\rightarrow$ higher Z, less pickup loading, slightly less drive.\\
Cin   & 0.1\,$\mu$F    & Input coupling cap; with the input impedance sets the low corner (38\,Hz). \emph{Smaller} $\rightarrow$ tighter, thinner low end.\\
Q1    & NPN (QNPN)     & First common-emitter gain stage; sets the drive into the clipper.\\
Rc1   & 10\,k$\Omega$  & Q1 collector load; sets first-stage gain. \emph{Larger} $\rightarrow$ more gain, more drive.\\
Ce1   & 47\,$\mu$F     & Emitter bypass; sets the AC gain. \emph{Smaller} $\rightarrow$ less gain, cleaner.\\
\sfgrouprow{3}{Clipper}\\
Ccp   & 0.22\,$\mu$F   & AC-couples into the clipper.\\
Rcp   & 47\,k$\Omega$  & Clipper load; sets how hard the diodes are driven. \emph{Smaller} $\rightarrow$ more drive, more distortion.\\
D1,\,D2 & Si (DSI)     & Symmetric silicon clipper. Lower-drop diodes (Ge/Schottky) $\rightarrow$ earlier, softer clip.\\
\sfgrouprow{3}{Recovery (Q2) \& output}\\
Q2    & NPN (QNPN)     & Recovery/output stage; restores level after the clipper.\\
Rc2   & 10\,k$\Omega$  & Q2 collector load; sets output level.\\
Cout  & 0.1\,$\mu$F    & Output coupling cap; keep $\geq$0.1\,$\mu$F or the low end thins under load.\\
Rtone & 25\,k$\Omega$  & Tone pot (treble-cut). \emph{More cut} $\rightarrow$ darker.\\
Ctone & 10\,nF         & Tone shunt cap; sets the treble-cut pivot.\\
Rvol  & 100\,k$\Omega$ & Master volume (modelled at maximum).\\
\end{tabular}
\caption{Complete bill of materials with per-part design notes.}
\end{table}

\begin{callout}{blue}{Pick the diodes by ear}
The clipper diodes dominate the character. \textbf{1N4148} silicon (modelled here) gives the classic
mid-forward drive; swapping to germanium or Schottky lowers the forward drop for an earlier, softer,
quieter clip, and an \emph{asymmetric} pair (one diode one way, two the other) adds even harmonics for
a grittier, more vocal break-up. The spec endpoints hold for any of these within the modelled range.
\end{callout}

\section{Build notes \& caveats}
These are the gaps between the verified \emph{model} and a physical \emph{build}. None affect the
10/10 simulation result; they are what to do at the bench.
\begin{itemize}
  \item \textbf{Use a real volume taper.} The volume is modelled at maximum; use an audio-taper
  100\,k$\Omega$ pot so the level sweep feels even.
  \item \textbf{Match the clipper diodes.} For true symmetry, match \texttt{D1} and \texttt{D2} forward
  drops; a few millivolts of mismatch is audible as a faint even-harmonic tint (often desirable).
  \item \textbf{Decouple the supply.} Add a reservoir cap across the 9\,V rail at the board; the model
  assumes a stiff supply.
  \item \textbf{The tone is a treble-cut, not a mid-scoop.} Voice and label it as a tone/darkness knob.
\end{itemize}

\section{Appendix --- netlist \& measurement methods}
\subsection{How the trickier numbers were measured}
\begin{itemize}
  \item \textbf{Distortion (thd\_max, thd\_clean)} --- two transient runs at 50\,mV and 5\,mV, each
  linearised and passed through ngspice's harmonic analysis; the THD is computed from the harmonic
  magnitudes and printed under two distinct names.
  \item \textbf{Symmetry (clip\_asym)} --- the positive and negative peak excursions of the clipped
  output are compared; the asymmetry is reported as a percentage of the full swing.
  \item \textbf{Tone range (tone\_db)} --- the 1\,kHz gain is measured at both tone-pot extremes and the
  difference is reported.
  \item \textbf{Input impedance (zin)} --- the small-signal input current is measured under a 1\,V AC
  drive and inverted.
\end{itemize}

\end{document}
